% Section 1: Introduction
% Source: Original Chinese draft - Section 1
% Status: PENDING TRANSLATION

Cancer remains one of the leading causes of mortality worldwide, significantly impacting both life expectancy and quality of life. Despite continuous advances in surgery, radiotherapy, chemotherapy, molecular targeted therapy, and immunotherapy, clinical benefits are often compromised by acquired drug resistance, systemic toxicity, and heterogeneity both within and across tumor lesions.

RNA-based therapeutics have garnered sustained attention due to their ability to directly modulate gene expression without relying on classical small-molecule binding pockets, offering distinct advantages for targeting "undruggable" proteins and complex signaling networks. However, the primary bottleneck for RNA drugs is not \textit{in vitro} efficacy, but rather their stability, biodistribution, and cellular uptake \textit{in vivo}. Naked RNA is susceptible to nuclease degradation, exhibits short circulation half-life, limited cellular internalization, and may trigger unintended immune responses. In essence, the delivery system largely defines the ceiling of RNA therapeutics.

Compared to fully synthetic nanomaterials, cell-secreted extracellular vesicles (EVs) possess natural membrane architecture, favorable biocompatibility, and intrinsic molecular transport capabilities, making them particularly promising as delivery chassis for RNA cargo. Plant-derived extracellular vesicles are emerging as an attractive alternative to mammalian EVs in this context. They offer more abundant source materials, generally lower production costs, and greater compatibility with edible plant resources and scalable extraction processes. Most studies further suggest that these vesicles exhibit relatively low immunogenicity and a favorable safety profile. Perhaps most compellingly, certain edible plant-derived vesicles retain stability in the gastrointestinal environment, presenting realistic possibilities for developing oral delivery systems.

Moreover, plant extracellular vesicles can carry small RNAs and participate in cross-species molecular regulation, positioning them not merely as "natural nano-shells" but as biologically active entities that provide new entry points for RNA-based cancer therapy. Based on these characteristics, plant-derived extracellular vesicles can be understood as a key transitional carrier bridging natural biological systems and engineered RNA delivery platforms.

Despite rapid progress in this field, several fundamental questions remain regarding RNA composition, mechanisms of cross-kingdom action, \textit{in vivo} biodistribution patterns, and scalable engineering pathways. This review systematically discusses recent advances in plant-derived extracellular vesicle-mediated RNA delivery for cancer therapy, with particular emphasis on their characteristics, mechanisms, engineering strategies, and applications across different tumor models, aiming to provide a clearer reference framework for future research and clinical translation.
